\subsection{Clipped objective and reference-policy KL}
\label{subsec:grpo-clipped-objective}

For sampled token $o_{i,t}$, recompute its log-probability under the current policy and
form the per-token ratio
\begin{equation}
  \rho_{i,t}(\theta)
  =\exp\!\left(
      \log\pi_\theta(o_{i,t}\mid q,o_{i,<t})
      -\log\pi_{\mathrm{old}}(o_{i,t}\mid q,o_{i,<t})
    \right).
  \label{eq:grpo-token-ratio}
\end{equation}
If one buffer is optimized repeatedly, an unconstrained $\rho_{i,t}A_i$ term can keep
rewarding movement even after the current policy is substantially different from the
policy that supplied the sample. GRPO limits that incentive with the clipped surrogate
\begin{equation}
  S_{i,t}(\theta)
  =\min\!\left(
      \rho_{i,t}A_i,
      \operatorname{clip}(\rho_{i,t},1-\epsilon,1+\epsilon)A_i
    \right).
  \label{eq:grpo-clipped-surrogate}
\end{equation}

Suppose $\epsilon=0.2$. For a positive advantage $A_i=1.5$ and ratio $\rho=1.5$, the raw
term is $2.25$, while the clipped candidate is $1.2(1.5)=1.8$; the minimum selects $1.8$.
The sampled action receives no additional credit for becoming more than $20\%$ more likely.
For a negative advantage $A_i=-1$ and ratio $\rho=0.5$, the raw term is $-0.5$ and the
clipped candidate is $-0.8$; the minimum selects $-0.8$. Reducing the probability farther
than the lower boundary earns no additional improvement. The sign-sensitive minimum caps
excess movement in the desired direction but does not hide movement in the wrong direction.

Clipping controls movement relative to the behavior policy for the current buffer. A
separate frozen reference policy limits cumulative drift from the model that began the RL
stage. Using the sampled-action KL estimator common in this objective family, define
\begin{equation}
  d_{i,t}
  =\log\pi_{\mathrm{ref}}(o_{i,t}\mid q,o_{i,<t})
   -\log\pi_\theta(o_{i,t}\mid q,o_{i,<t}),
  \qquad
  \widehat D_{i,t}=\exp(d_{i,t})-d_{i,t}-1.
  \label{eq:grpo-sampled-kl}
\end{equation}
The token objective is
\begin{equation}
  J_{i,t}^{\mathrm{GRPO}}(\theta)
  =S_{i,t}(\theta)-\beta\widehat D_{i,t}.
  \label{eq:grpo-token-objective}
\end{equation}
This token objective now contains both pieces promised by the top-level formula. Averaging
it first over tokens and then equally over responses reassembles the complete loss:
\begin{equation}
  \mathcal{L}_{\mathrm{GRPO}}(\theta)
  =-\frac{1}{G}\sum_{i=1}^{G}\frac{1}{T_i}\sum_{t=1}^{T_i}
    J_{i,t}^{\mathrm{GRPO}}(\theta)
  =-\frac{1}{G}\sum_{i=1}^{G}\frac{1}{T_i}\sum_{t=1}^{T_i}
    \left(S_{i,t}(\theta)-\beta\widehat D_{i,t}\right).
  \label{eq:grpo-reassembled-loss}
\end{equation}
Equation~\ref{eq:grpo-reassembled-loss} is Equation~\ref{eq:grpo-full-loss} written with
the shorthand defined in this subsection. The clipping coefficient $\epsilon$ controls
short-term reuse of one behavior batch; the KL coefficient $\beta$ controls long-term
movement from the reference model.
Historically, this clipped probability-ratio construction comes from Proximal Policy
Optimization; GRPO adopts it after replacing the learned value baseline with its grouped
reward baseline \citep{schulman2017ppo}.
